%%%% ABSTRACT
%%
%% Versão do resumo para idioma de divulgação internacional.

\begin{abstractutfpr}
This Final Project will present the process of refactoring, modernizing, and expanding the E-Museu web system, a digital platform designed for cataloging and disseminating items from the history of computing. The justification for carrying out this work will be based on the technical limitations identified in the current version of the application, which presents maintenance difficulties, lack of code standardization, insufficient automated tests, and the absence of essential functionalities for the continuity of the project. In this context, the work will aim to improve the system’s internal structure, increase its long-term stability and sustainability, and simultaneously implement new features that will expand the platform’s practical usefulness. As methodology, an approach based on software engineering best practices will be adopted, including the application of design patterns, module reorganization, versioning standardization, and the configuration of development, testing, and production environments. Additionally, refactoring principles will be applied to improve code readability, modularity, and performance, while automated tests will be incorporated to ensure the reliability of features and prevent future issues. Among the expected results are the complete restructuring of the codebase, significant improvements in project organization and traceability, the implementation of security mechanisms, increased media upload capacity, system internationalization, and the generation of labels for the physical cataloging of items. It is expected that the conclusion will demonstrate that the set of improvements applied will make the E-Museu more robust and aligned with modern development practices, establishing a solid foundation for future evolutions and contributing to the preservation of technological memory.
\end{abstractutfpr}